\chapter{Le processus de recherche}\label{ch:rs:the-research-process}

Un chercheur présente un backtest le vendredi après-midi~: un signal de retournement sur des actions liquides, deux ans de données quotidiennes, un
ratio de Sharpe annuel de 2.4 après coûts. Le lundi, le responsable de la recherche pose deux questions. Combien de variantes ont été testées avant
celle-ci~? Lequel est venu en premier, la raison pour laquelle le signal devrait fonctionner ou le graphique qui affirme qu'il fonctionne~? Personne n'a noté ni l'un ni l'autre, et
sans elles, le chiffre affiché à l'écran ne peut pas être jugé. Ce livre porte sur le métier qui transforme les données de marché en stratégies auxquelles une
maison peut se fier~: les prédicteurs et la façon de les mesurer, les prévisions, les backtests à quatre niveaux de fidélité, la mesure de la performance et du risque
et la construction de portefeuilles. Il commence par le processus, car tout outil ultérieur peut être mis en échec par un
processus qui n'enregistre pas ce qu'il a essayé.

\section{L'hypothèse d'abord}

\begin{definition}[Hypothèse de recherche, justification économique]\label{def:rs:the-research-process:hypothesis}
Une \emph{hypothèse de recherche}\index{hypothèse de recherche} (research hypothesis) est un énoncé réfutable, formulé avant le test, portant sur un
mécanisme qui fait bouger les prix~: quels instruments, à quel horizon, quelle grandeur observable prédit quoi, et quel résultat du
test la réfuterait. Sa \emph{justification économique}\index{justification économique} (economic rationale) est la réponse à la question «~qui paie,
et pourquoi~?~»~: un risque que la stratégie porte et que d'autres paieront pour s'en décharger, un service qu'elle rend (liquidité, immédiateté, le
portage de l'inventaire), une contrainte qui force d'autres acteurs à négocier à un mauvais prix (un rééquilibrage d'indice, un appel de marge, un
mandat), ou une information traitée plus vite ou mieux que ne le reflète déjà le cours.
\end{definition}

La justification compte pour deux raisons. Elle indique au chercheur où regarder, et réduit ainsi l'espace de recherche~; et elle prédit
davantage que le résultat principal, ce qui offre à l'hypothèse des moyens d'être réfutée. Un scénario de fourniture de liquidité appliqué aux retournements
à court terme prédit que le retournement est plus fort dans les actions moins liquides, plus fort les jours où la liquidité est rare, et
plus faible après la baisse du coût de l'immédiateté sur le marché. Une stratégie dont le backtest réussit alors que chacune des prédictions secondaires échoue est
vraisemblablement en train d'ajuster autre chose.

\begin{example}[Une hypothèse consignée par écrit]\label{ex:rs:the-research-process:written}
\emph{«~Les actions qui chutent davantage que leur secteur sur une journée, sans nouvelle particulière, reprennent une partie du mouvement au cours des cinq
jours suivants, parce que les vendeurs exigent de l'immédiateté et que les teneurs de marché sont rémunérés pour la fournir. L'effet devrait être plus fort
dans la moitié inférieure de l'univers en valeur traitée et dans le décile supérieur de volatilité de marché, et absent les jours où
une publication de résultats. Si le coefficient d'information à cinq jours du rendement d'une journée, relatif au secteur, n'est pas négatif dans les deux
moitiés de l'échantillon, l'hypothèse est rejetée.~»} Chacune de ses clauses peut être testée, et la dernière phrase dit
d'avance à quoi ressemble un échec.
\end{example}

Le cas inverse est instructif. Arnott, Harvey et Markowitz (2019) ouvrent leur protocole de backtest par une stratégie long--short
sur des actions du NYSE, mise au point sur 1963--1988 et «~validée hors échantillon avec des résultats encore plus forts~» sur
1989--2015~: un alpha proche de 6~\% par an, aucune corrélation significative avec le marché ni avec les facteurs bien connus, un taux de rotation inférieur à
10~\% par an. Elle a été construite à partir des lettres des symboles boursiers, choisis parmi des milliers de combinaisons de lettres. Un
test hors échantillon ne l'a pas sauvée, car rien d'autre que la recherche ne la soutenait.

Ce que dit un test réussi dépend de la proportion d'hypothèses vraies parmi celles qu'une équipe teste, un chiffre inconfortable que
Ioannidis (2005) a placé au centre d'un débat sur la recherche publiée.

\begin{proposition}[Ce que vaut un backtest significatif]\label{prop:rs:the-research-process:ppv}
Supposons qu'une fraction $\pi$ des hypothèses testées par une équipe soient vraies, et que chacune soit testée par $k$ étapes indépendantes, chacune avec
une taille $\alpha$ et une puissance $1 - \beta$ face aux hypothèses vraies. La probabilité qu'une hypothèse franchissant les $k$ étapes soit vraie
(sa valeur prédictive positive) est
\[
\mathrm{PPV}_k = \frac{\pi(1 - \beta)^k}{\pi(1 - \beta)^k + (1 - \pi)\alpha^k}.
\]
\end{proposition}

\begin{proof}
D'après la formule de Bayes, avec $\P(\text{réussite} \mid \text{vraie}) = (1 - \beta)^k$ et $\P(\text{réussite} \mid \text{fausse}) = \alpha^k$ par
indépendance des étapes.
\end{proof}

Avec une étape, une probabilité a priori de 5~\%, la taille conventionnelle de 5~\% et une puissance d'un demi (un edge modeste sur quelques années de données),
$\mathrm{PPV}_1 = 0.345$~: deux backtests sur trois qui réussissent sont faux. Une deuxième étape indépendante (un \omterm{def:rs:the-research-process:prereg}{échantillon réservé} (holdout), un autre marché)
le porte à 0.840 et la troisième étape à 0.981 (la \cref{fig:rs:the-research-process:ppv}). Un meilleur test aide moins qu'un meilleur a priori~:
avec une probabilité a priori de 10~\% et une puissance de 80~\%, une seule étape donne déjà 0.64.

\begin{omfigure}
\begin{tikzpicture}
\begin{semilogxaxis}[width=10cm,height=5.2cm,xlabel={probabilité a priori $\pi$ qu'une idée testée soit vraie},
  ylabel={$\mathrm{PPV}_k$},tick label style={font=\small},label style={font=\small},xmin=0.002,xmax=0.5,ymin=0,ymax=1,
  grid=major,grid style={gray!25},legend columns=3,legend style={font=\footnotesize,at={(0.5,-0.32)},anchor=north,
  draw=none,fill=none,/tikz/every even column/.append style={column sep=8pt}},table/col sep=comma]
\addplot[omThm,very thick,dashed] table[x=prior,y=one]{figdata/research/01-the-research-process/ppv.csv};
\addplot[omDef,very thick] table[x=prior,y=two]{figdata/research/01-the-research-process/ppv.csv};
\addplot[black!70,very thick,dotted] table[x=prior,y=three]{figdata/research/01-the-research-process/ppv.csv};
\addplot[black!50,thin,sharp plot] coordinates {(0.05,0) (0.05,1)};
\legend{une étape,deux étapes,trois étapes}
\end{semilogxaxis}
\end{tikzpicture}

\omcaption{Probabilité qu'une idée ayant franchi $k$ étapes indépendantes soit vraie, selon la part d'idées vraies parmi celles
testées, chaque étape ayant une taille de 5~\% et une puissance de 50~\%. Pour une probabilité a priori de 5~\% (ligne verticale)~: 0.345, 0.840 et 0.981. Données~: le
module du chapitre, \cref{prop:rs:the-research-process:ppv}.}\label{fig:rs:the-research-process:ppv}
\end{omfigure}

\section{Le journal de recherche et le nombre d'essais}

\begin{definition}[Journal de recherche, nombre d'essais]\label{def:rs:the-research-process:log}
Un \emph{journal de recherche}\index{journal de recherche} (research log) est un registre en ajout seul, tenu au fil du travail, de chaque hypothèse enregistrée
et de chaque test exécuté~: ses paramètres, le code et les données utilisés, ses résultats et la date de son exécution. Le
\emph{nombre d'essais}\index{nombre d'essais} (trial count) d'un résultat est le nombre de tests portant sur la même famille d'hypothèses consignés dans le journal avant le choix du
résultat, y compris les variantes qui ont échoué et ont été écartées.
\end{definition}

Le \omterm{def:rs:the-research-process:log}{nombre d'essais} est ce dont a besoin le livre~4, chapitre~12, pour corriger un résultat de la recherche qui l'a rendu possible~: ses corrections pour tests
multiples et son ratio de Sharpe déflaté prennent toutes deux le \omterm{def:rs:the-research-process:log}{nombre d'essais} comme donnée d'entrée, et un chercheur qui ne le consigne pas
ne peut pas le fournir. Les essais les plus difficiles à compter sont ceux qui ne sont jamais exécutés, seulement envisagés~: le jardin aux sentiers qui bifurquent
de ce chapitre. Le journal ne les capte pas, mais il capte tout le reste et rend la différence visible.

\begin{proposition}[La meilleure parmi de nombreuses variantes sans valeur]\label{prop:rs:the-research-process:best}
Soient $n$ variantes sans edge, chacune testée en backtest sur $Y$ années de rendements quotidiens, de sorte que chaque ratio de Sharpe annuel estimé soit
approximativement $\mathcal N(0, 1/Y)$, et soient les estimations équicorrélées de corrélation $\rho$. Alors
\[
\P\Bigl(\max_i \widehat{\mathrm{SR}}_i \ge c\Bigr) = \E_W\Bigl[1 - \Phi\Bigl(\frac{c\sqrt Y - \sqrt\rho\,W}{\sqrt{1 - \rho}}\Bigr)^n\Bigr],
\qquad W \sim \mathcal N(0, 1),
\]
ce qui vaut $1 - \Phi(c\sqrt Y)^n$ pour des variantes indépendantes.
\end{proposition}

\begin{proof}
Écrivons les estimations standardisées sous la forme $Z_i = \sqrt\rho\,W + \sqrt{1 - \rho}\,\varepsilon_i$, avec des normales centrées réduites indépendantes~; sachant
$W$, les $Z_i$ sont indépendants, donc $\P(\max_i Z_i < c\sqrt Y \mid W)$ est le produit des $n$ probabilités conditionnelles.
Intégrer par rapport à $W$.
\end{proof}

Les chiffres sont saisissants (\cref{fig:rs:the-research-process:best}). Une seule stratégie sans valeur, backtestée sur deux ans, affiche
un ratio de Sharpe de 2 ou plus avec une probabilité de 0.23~\%. La meilleure parmi vingt indépendantes le fait dans 4.6~\% des cas, et la
meilleure parmi mille dans 90~\% des cas~: le meilleur parmi mille atteint en moyenne 2.30, au-dessus du ratio qui a impressionné la salle
le vendredi. Un an de données fait d'un ratio de Sharpe de 2 une quasi-certitude pour le meilleur parmi mille~; cinq ans ramènent la
probabilité à 0.39~\%. La corrélation entre les variantes aide, car des variantes corrélées sont moins de chances indépendantes~: vingt variantes
corrélées à 0.6 atteignent 2 avec une probabilité de 2.6~\%.

\begin{omfigure}
\begin{tikzpicture}
\begin{semilogxaxis}[width=10cm,height=5.4cm,xlabel={nombre de variantes sans valeur testées $n$},
  ylabel={$\P(\text{meilleur Sharpe} \ge 2)$},tick label style={font=\small},label style={font=\small},xmin=1,xmax=10000,
  ymin=0,ymax=1,grid=major,grid style={gray!25},legend columns=2,legend cell align=left,legend style={font=\footnotesize,
  at={(0.5,-0.32)},anchor=north,draw=none,fill=none,/tikz/every even column/.append style={column sep=14pt}},
  table/col sep=comma]
\addplot[omThm,very thick,dashed] table[x=n,y=y1]{figdata/research/01-the-research-process/best.csv};
\addplot[omDef,very thick] table[x=n,y=y2]{figdata/research/01-the-research-process/best.csv};
\addplot[black!70,very thick,dotted] table[x=n,y=y5]{figdata/research/01-the-research-process/best.csv};
\addplot[omDef!60,thick,dashdotted] table[x=n,y=y2corr]{figdata/research/01-the-research-process/best.csv};
\legend{1 an de données,2 ans,5 ans,{2 ans, corrélation 0.6}}
\end{semilogxaxis}
\end{tikzpicture}

\omcaption{Probabilité que le meilleur de $n$ stratégies sans avantage affiche un ratio de Sharpe annuel d'au moins 2, selon la durée
du backtest~; la courbe en tirets-points correspond à des estimations des variantes corrélées à 0.6. Avec deux ans de données~: 0.23~\% pour
une, 4.6~\% pour vingt, 90~\% pour mille variantes indépendantes. Données~: le module du chapitre,
\cref{prop:rs:the-research-process:best}.}\label{fig:rs:the-research-process:best}
\end{omfigure}

\begin{remark}[Nombre de combinaisons]\label{rem:rs:the-research-process:combinations}
Les essais se multiplient plus vite que ne le remarquent les chercheurs. Vingt variables candidates avec toutes leurs interactions deux à deux ne
donnent pas 22 tests mais $20 + \binom{20}{2} = 210$~; Arnott, Harvey et Markowitz font le même décompte. Un changement d'univers, de jour de
rééquilibrage ou d'hypothèse de coûts après avoir vu un résultat est aussi un essai, et doit figurer dans le \omterm{def:rs:the-research-process:log}{journal de recherche}.
\end{remark}

\section{Ce qui compte comme preuve}

\begin{definition}[Pré-enregistrement, échantillon réservé]\label{def:rs:the-research-process:prereg}
Le \emph{pré-enregistrement}\index{pré-enregistrement} consiste à consigner une hypothèse et son plan d'analyse (données, échantillon,
Un \emph{échantillon réservé}\index{echantillon@échantillon réservé} (holdout set) est une partie de
l'ensemble de données mise de côté avant le début de la recherche et utilisée une seule fois, pour le test final d'une analyse pré-enregistrée.
\end{definition}

Nosek et ses coauteurs (2018) résument l'enjeu en une distinction~: le \omterm{def:rs:the-research-process:prereg}{pré-enregistrement} sépare les analyses qui testent des prédictions de celles
qui les engendrent a posteriori, et la crédibilité d'un résultat dépend de celle à laquelle il appartient. En finance, la distinction est plus
difficile à maintenir que dans un essai clinique, car les données existent avant la question. Un \omterm{def:rs:the-research-process:prereg}{échantillon réservé} permet de la retrouver,
à condition qu'il soit réellement tenu à l'écart.

\begin{method}[Évaluer les preuves en faveur d'une stratégie]\label{met:rs:the-research-process:evidence}
Classer ce que la recherche a produit, du plus faible au plus fort, et exiger au moins trois niveaux~:
\begin{enumerate}
\item un ajustement dans l'échantillon, aussi bon soit-il (il mesure la recherche autant que l'idée)~;
\item le même test sur des données postérieures qui n'ont servi à rien choisir~;
\item les prédictions secondaires du mécanisme~: d'autres marchés, d'autres périodes, les sous-ensembles où la \omterm{def:rs:the-research-process:hypothesis}{justification économique} veut que
l'effet soit plus fort ou absent~;
\item un test pré-enregistré sur l'\omterm{def:rs:the-research-process:prereg}{échantillon réservé}, exécuté une seule fois~;
\item un trading en direct de faible taille, mesuré par rapport au backtest des mêmes jours.
\end{enumerate}
Chaque niveau ne compte que s'il était indépendant des choix faits aux niveaux précédents.
\end{method}

Un \omterm{def:rs:the-research-process:prereg}{échantillon réservé} qui a déjà été consulté n'est pas indépendant. Si un chercheur le vérifie cinq fois pendant le réglage, à chaque fois au seuil de 5~\%,
la probabilité qu'une stratégie sans valeur passe au moins une fois est $1 - 0.95^5 = 22.6~\%$, et non 5~\%. Dans la
\cref{prop:rs:the-research-process:ppv}, la seconde étape a alors une taille de 0.226 au lieu de 0.05, et avec un a priori de 5~\%, les deux
étapes donnent ensemble une VPP de 0.54 au lieu de 0.84. Un \omterm{def:rs:the-research-process:prereg}{échantillon réservé} est dépensé dès qu'il influence un choix.

\section{De l'idée à la production}

\begin{definition}[Jalon de validation, critère d'abandon, revue de recherche]\label{def:rs:the-research-process:gate}
Un \emph{jalon de validation}\index{jalon de validation} (stage gate) est un point de décision entre deux étapes de la vie d'une stratégie, auquel elle doit satisfaire à des
critères fixés à l'avance pour continuer. Un \emph{critère d'abandon}\index{critère d'abandon} (kill criterion) est une règle, écrite avant le début d'une étape,
qui arrête la stratégie lorsque les preuves qu'elle produit tombent sous un seuil. Une \emph{revue de recherche}\index{revue de recherche} (research review)
est l'examen d'un résultat par des personnes qui ne l'ont pas produit, avec accès au \omterm{def:rs:the-research-process:log}{journal de recherche}.
\end{definition}

\begin{omfigure}
\begin{tikzpicture}[font=\footnotesize,
  stg/.style={draw=omDef,fill=omDef!8,rounded corners=2pt,text width=1.45cm,minimum height=0.95cm,align=center,font=\scriptsize},
  gt/.style={->,>=Stealth,thick,omDef}, kl/.style={->,>=Stealth,thick,omThm}]
\node[stg] (a) at (0,0) {idée et\\ hypothèse};
\node[stg] (b) at (2.10,0) {exploration};
\node[stg] (c) at (4.20,0) {test sur\\ échantillon réservé};
\node[stg] (d) at (6.30,0) {revue};
\node[stg] (e) at (8.40,0) {trading fantôme\\ (shadow trading)};
\node[stg] (f) at (10.50,0) {petit\\ direct};
\node[stg] (g) at (12.60,0) {montée en charge};
\draw[gt] (a) -- (b); \draw[gt] (b) -- (c); \draw[gt] (c) -- (d); \draw[gt] (d) -- (e); \draw[gt] (e) -- (f); \draw[gt] (f) -- (g);
\node[draw=omThm,fill=omThm!8,rounded corners=2pt,minimum width=10cm,minimum height=0.7cm] (k) at (6.3,-1.9)
  {arrêtée~: motif consigné dans le journal de recherche};
\foreach \s in {b,c,d,e,f} \draw[kl] (\s.south) -- (\s.south |- k.north);
\node[omExo,align=center] at (6.3,1.05) {chaque flèche est un jalon de validation aux critères fixés à l'avance};
\end{tikzpicture}

\omcaption{La vie d'une stratégie. Chaque étape est atteinte par un jalon dont les critères ont été fixés avant le début de l'étape, et
chacune peut se terminer par un arrêt dont le motif est consigné~: les arrêts sont des essais, et le journal les compte.}\label{fig:rs:the-research-process:stages}
\end{omfigure}

Le pipeline de la \cref{fig:rs:the-research-process:stages} est générique~; les entreprises nomment les étapes différemment et certaines les fusionnent.
Sa logique ne change pas. L'exploration peut chercher librement, car elle est consignée. Le test sur l'\omterm{def:rs:the-research-process:prereg}{échantillon réservé} exécute une seule analyse pré-enregistrée
sur celui-ci. La revue lit le journal en même temps que le résultat. Le trading fantôme exécute la stratégie sur des données en direct, sans ordres
(chapitre~19), et le trading en direct de faible taille mesure ce que le simulateur ne peut pas mesurer~: les exécutions, l'impact et la discipline des personnes
qui la font tourner. Un \omterm{def:rs:the-research-process:gate}{critère d'abandon} demande le même soin que les critères d'entrée, car les historiques en direct sont courts et bruités.

\begin{proposition}[Ce qu'un historique en direct dit contre son backtest]\label{prop:rs:the-research-process:kill}
Si le ratio de Sharpe annuel véritable d'une stratégie était le $S$ issu du backtest, la probabilité que $Y$ années de rendements quotidiens en direct affichent un
ratio de Sharpe annuel de $s$ ou moins vaut approximativement $\Phi\bigl((s - S)\sqrt Y\bigr)$. Distinguer un $S$ véritable de zéro par un
test unilatéral de taille $\alpha$ et de puissance $1 - \beta$ demande environ $Y = \bigl((z_{1-\alpha} + z_{1-\beta})/S\bigr)^2$ années.
\end{proposition}

\begin{proof}
L'estimation est approximativement $\mathcal N(S, 1/Y)$~; la seconde affirmation relève du livre~4, chapitre~12, avec les rôles écrits pour un
historique en direct.
\end{proof}

Une stratégie testée en backtest à 1.5 qui rapporte un ratio de Sharpe de 0.2 lors de sa première année en direct est décevante mais pas décisive~:
selon la prétention même du backtest, une année aussi mauvaise a une probabilité de 9.7~\%. Un historique en direct qui doit séparer 1.5 de zéro avec une taille de 5~\%
et une puissance de 80~\% demande 2.7 années. D'où la règle de pratique~: un \omterm{def:rs:the-research-process:gate}{critère d'abandon} est un seuil sur une statistique fixée
à l'avance, assorti d'une durée minimale, et souvent d'un second déclencheur portant sur une grandeur mesurée plus vite que les rendements, comme
les exécutions face au simulateur ou les coûts réalisés.

\section{La recherche comme portefeuille de paris}

Un groupe de recherche est lui-même une stratégie~: il consacre du temps de chercheur et des données à des idées à faible probabilité a priori et à gain asymétrique. Son
production attendue sur un an découle de la \cref{prop:rs:the-research-process:ppv}. Un groupe qui teste 200 idées par an, dont
5~\% sont réelles, avec des étapes de taille 5~\% et de puissance 50~\%, laisse passer 5 idées réelles et 9.5 fausses lors de l'exploration, puis 2.5 réelles
et 0.475 fausse lors d'une seconde étape indépendante. Le premier jalon divise par deux les idées réelles et divise les fausses par vingt~;
ce compromis en vaut presque toujours la peine, car une stratégie fausse qui atteint la production coûte ses pertes de trading, le
capital qu'elle immobilise et les mois qui passent avant que le \omterm{def:rs:the-research-process:gate}{critère d'abandon} ne se déclenche.

Trois règles empiriques en découlent. Relever la probabilité a priori avant de relever la puissance~: une hypothèse qui repose sur une
\omterm{def:rs:the-research-process:hypothesis}{justification économique} et sur des prédictions secondaires vaut plus qu'une recherche plus large. Garder les étapes indépendantes~: un \omterm{def:rs:the-research-process:prereg}{échantillon réservé} utilisé une seule fois, un marché non examiné, une période non étudiée.
Et compter~: le \omterm{def:rs:the-research-process:log}{journal de recherche} ne coûte presque rien, et c'est le seul document à partir duquel quiconque peut plus tard distinguer une découverte d'une simple exploration.

\section{Tutoriel~: consigner une recherche}\label{tut:rs:the-research-process:log}

\begin{tutorial}
\textbf{Objectif.} Lancer une recherche de cinquante variantes sur des données où rien ne fonctionne, avec chaque essai inscrit dans le \omterm{def:rs:the-research-process:log}{journal de recherche}, et laisser le journal
déflater le meilleur résultat. \textbf{État final~:} la \cref{fig:rs:the-research-process:search}~; le meilleur résultat, avec un lookback de deux jours, affiche un ratio de Sharpe annuel de
2.14, un ratio de Sharpe probabiliste de 0.999 comme test isolé et un ratio de Sharpe déflaté de 0.77 compte tenu des
cinquante essais consignés.
\begin{tutsteps}
\item \textbf{Le journal.} L'empreinte de chaque entrée couvre son contenu et l'empreinte précédente, de sorte que modifier une entrée passée rompt
chaque lien suivant~; \texttt{verify()} trouve le premier lien rompu.
\omcode{code/firm/researchlog/firm_researchlog.py}{44}{67}{Entrées chaînées par hachage du journal de recherche.}\label{lst:rs:the-research-process:chain}
\item \textbf{La recherche.} Enregistrer d'abord l'hypothèse, puis consigner cinquante règles de retournement (conserver moins le signe du
rendement passé de $L$ jours, $L = 1, \dots, 50$) sur deux ans d'une marche aléatoire sans dérive, de volatilité quotidienne de 1~\%.
\omcode{code/research/01-the-research-process/python/rs_process.py}{96}{117}{Cinquante essais consignés d'une règle de retournement sur du bruit.}\label{lst:rs:the-research-process:search}
\item \textbf{Déflater.} Appeler \texttt{log.deflated(family, sr, n\_obs)} sur la famille \texttt{reversal} avec le ratio de Sharpe quotidien du meilleur résultat~: la référence est
l'espérance du meilleur parmi le nombre consigné d'essais nuls. Puis exécuter \texttt{fig\_process.py}.
\end{tutsteps}
\textbf{Pour aller plus loin.} Ajouter une seconde famille (règles de momentum) au même journal et vérifier que chaque famille est déflatée par
son propre \omterm{def:rs:the-research-process:log}{nombre d'essais}~; modifier une métrique consignée dans le fichier JSON et observer où \texttt{verify()} la signale.
\end{tutorial}

\begin{omfigure}
\begin{tikzpicture}
\begin{axis}[width=6.6cm,height=5cm,xlabel={lookback $L$ (jours)},ylabel={ratio de Sharpe annuel},
  tick label style={font=\small},label style={font=\small},xmin=0,xmax=51,ymin=-2,ymax=2.5,grid=major,grid style={gray!25},
  table/col sep=comma]
\addplot[omDef,thick,mark=*,mark size=1pt] table[x=lookback,y=sr]{figdata/research/01-the-research-process/search.csv};
\addplot[omThm,only marks,mark=*,mark size=2.4pt] coordinates {(2,2.141)};
\end{axis}
\end{tikzpicture}\hspace{4mm}%
\begin{tikzpicture}
\begin{semilogxaxis}[width=6.6cm,height=5cm,xlabel={essais consignés $N$},ylabel={ratio de Sharpe déflaté},
  tick label style={font=\small},label style={font=\small},xmin=1,xmax=1000,ymin=0,ymax=1.05,grid=major,grid style={gray!25},
  table/col sep=comma]
\addplot[omDef,very thick] table[x=n,y=dsr]{figdata/research/01-the-research-process/deflate.csv};
\addplot[black!50,thin,sharp plot] coordinates {(50,0) (50,1.05)};
\end{semilogxaxis}
\end{tikzpicture}

\omcaption{À gauche~: ratios de Sharpe annuels des cinquante règles de retournement sur une marche aléatoire simulée de deux ans~; le meilleur résultat (point) affiche
2.14 pour un lookback de deux jours, cinq règles dépassent 1 et la moyenne vaut 0.07. À droite~: ratio de Sharpe déflaté du meilleur résultat en fonction du
\omterm{def:rs:the-research-process:log}{nombre d'essais} que le journal enregistre~: 0.999 pour un seul, 0.77 pour les cinquante réellement effectués (ligne verticale). Données~: le tutoriel du chapitre,
avec germe fixé (la première histoire tirée).}\label{fig:rs:the-research-process:search}
\end{omfigure}

\section{Construction~: le journal de recherche}\label{bld:rs:the-research-process:researchlog}

\begin{build}
\bfield{Rôle} La mémoire de la petite entreprise de ce qu'a tenté sa recherche. Chaque outil ultérieur de ce livre qui exécute un test
(les backtesters de la partie~IV, les outils de surapprentissage du chapitre~20, le flux de travail du chapitre~29) inscrit ici un essai, et la
revue d'un résultat commence par la lecture de sa famille.
\bfield{Interface} \texttt{ResearchLog(path)}~; \texttt{register(family, hypothesis, ts)}~;
\texttt{trial(family, params, metrics, ts, code\_hash, data\_id)} renvoyant un \texttt{Entry}~; \texttt{trial\_count(family)}~;
\texttt{post\_hoc(family)}~; \texttt{metric(family, name)}~; \texttt{verify()}~; \texttt{deflated(family, sr, n\_obs, skew, kurt)}.
\bfield{Règles} Journal en ajout seul, en lignes JSON~; chaque entrée porte l'empreinte de son contenu et de celle de l'entrée précédente~; les horodatages sont
fournis par l'appelant, jamais lus sur l'horloge, de sorte qu'un journal rejoué soit identique~; une hypothèse par famille, et une famille
dont la première entrée est un essai est signalée comme post hoc~; la déflation utilise \texttt{firm.multitest} du livre~4.
\bfield{Tests d'acceptation} \texttt{code/firm/researchlog/tests/}~: une chaîne intacte se vérifie~; une métrique modifiée est localisée~; une famille
post hoc est signalée et une seconde hypothèse refusée~; le journal survit à un rechargement et continue de compter~; le ratio de Sharpe déflaté
baisse à mesure que le \omterm{def:rs:the-research-process:log}{nombre d'essais} croît.
\bfield{Pour aller plus loin} Signer les entrées avec une clé par chercheur~; compter le nombre effectif d'essais d'une famille à partir de la
corrélation de ses séries de P\&L consignées (livre~4, chapitre~12) au lieu du nombre brut.
\end{build}

\begin{omsources}
\item J.~P.~A.~Ioannidis, ``Why most published research findings are false'', \emph{PLoS Medicine} 2(8), e124, 2005.
\item R.~D.~Arnott, C.~R.~Harvey et H.~Markowitz, ``A backtesting protocol in the era of machine learning'', \emph{Journal of
Financial Data Science} 1(1), 2019.
\item B.~A.~Nosek, C.~R.~Ebersole, A.~C.~DeHaven et D.~T.~Mellor, ``The preregistration revolution'', \emph{Proceedings of the
National Academy of Sciences} 115(11), 2018.
\item D.~H.~Bailey, J.~M.~Borwein, M.~López de Prado et Q.~J.~Zhu, ``Pseudo-mathematics and financial charlatanism: the effects
of backtest overfitting on out-of-sample performance'', \emph{Notices of the American Mathematical Society} 61(5), 2014.
\item M.~López de Prado, ``The 10 reasons most machine learning funds fail'', \emph{Journal of Portfolio Management} 44(6), 2018.
\item C.~R.~Harvey, Y.~Liu et H.~Zhu, ``\dots and the cross-section of expected returns'', \emph{Review of Financial Studies}
29(1), 2016.
\end{omsources}

\section{Exercices}

\begin{exercise}[$\star$]\label{exo:rs:the-research-process:1}
Une équipe estime que 10~\% de ses idées sont réelles et teste chacune avec une puissance de 80~\% à une taille de 5~\%. Quelle fraction des idées qui passent
sont réelles~? Et si la probabilité a priori est de 2~\%~?
\end{exercise}

\begin{exercise}[$\star$]\label{exo:rs:the-research-process:2}
Vingt stratégies indépendantes sans valeur sont chacune testées unilatéralement à $t \ge 2$. Quelle est la probabilité qu'au moins une
passe, et combien en passent en moyenne~?
\end{exercise}

\begin{exercise}[$\star$]\label{exo:rs:the-research-process:3}
Une stratégie testée en backtest à un ratio de Sharpe annuel de 1.2 rapporte $-0.5$ au cours de ses six premiers mois en direct. Quelle est la probabilité d'un semestre
aussi mauvais si le backtest était juste~?
\end{exercise}

\begin{exercise}[$\star\star$]\label{exo:rs:the-research-process:4}
Combien d'années de rendements quotidiens en direct séparent un ratio de Sharpe annuel de 0.8 de zéro avec un test unilatéral de taille 5~\% et une puissance de
90~\%~? Comparer avec la réponse pour 1.5 à 80~\% de puissance dans le texte.
\end{exercise}

\begin{exercise}[$\star\star$]\label{exo:rs:the-research-process:5}
Un chercheur teste trente variables et tous leurs produits deux à deux. Combien de tests cela représente-t-il, et quelle statistique $t$ une
correction de Bonferroni unilatérale à 5~\% exige-t-elle~?
\end{exercise}

\begin{exercise}[$\star\star$]\label{exo:rs:the-research-process:6}
Une équipe teste 150 idées par an avec une probabilité a priori de 4~\% à travers trois étapes indépendantes de taille 5~\% et de puissance 60~\%. Combien de stratégies réelles et
combien de fausses stratégies atteignent la production chaque année, et quelle est la VPP à chaque étape~?
\end{exercise}

\begin{exercise}[$\star\star\star$]\label{exo:rs:the-research-process:7}
\emph{Programmation.} Avec \texttt{p\_best}, calculer la probabilité que la meilleure de mille variantes sans valeur, testées en backtest sur deux
ans, atteigne un ratio de Sharpe de 2 lorsque les variantes sont équicorrélées à 0.3, 0.6 et 0.9. Quelle corrélation rend mille
variantes à peu près aussi utiles que vingt variantes indépendantes~?
\end{exercise}

\begin{exercise}[$\star\star\star$]\label{exo:rs:the-research-process:8}
\emph{Trouver la faille.} «~Nous avons trouvé le signal sur 2015--2020 et l'avons confirmé sur l'\omterm{def:rs:the-research-process:prereg}{échantillon réservé} 2021--2023 avec un ratio de Sharpe de 1.4.
Nous avons ensuite ajusté deux paramètres sur l'\omterm{def:rs:the-research-process:prereg}{échantillon réservé}, ce qui l'a porté à 2.1. La stratégie a passé deux étapes indépendantes et nous
proposons de la trader.~»
\end{exercise}

\section{Problème~: Le backtest du vendredi}

\begin{problem}[{Problème du week-end --- à quelle fréquence une recherche sans valeur produit-elle un ratio de Sharpe de 2~?}]\label{pb:rs:the-research-process:1}
Un chercheur sans avantage essaie chaque semaine vingt variantes d'un signal, cinquante semaines par an. Chaque variante est testée en backtest sur deux
ans de données quotidiennes, les vingt variantes d'une semaine sont équicorrélées à 0.6, et les recherches de semaines différentes sont indépendantes. Un
backtest «~impressionne~» à un ratio de Sharpe annuel de 2.

\medskip\noindent\textbf{Partie I --- Un backtest.}
\begin{enumerate}
\item Quelle est la loi du ratio de Sharpe annuel estimé d'une variante~?
\item Quelle est la probabilité qu'une variante impressionne~?
\item Quelle est l'espérance du nombre de variantes impressionnantes en un an, en ignorant la corrélation~?
\item Pourquoi cette espérance ne répond-elle pas à la question que pose le responsable de la recherche~?
\end{enumerate}

\medskip\noindent\textbf{Partie II --- Une semaine, une année.}
\begin{enumerate}[resume]
\item Quelle est la probabilité que la meilleure des vingt variantes d'une semaine impressionne si elles étaient indépendantes~?
\item Et avec leur corrélation de 0.6~?
\item Quelle est la probabilité d'au moins un backtest impressionnant en un an~?
\item Que vaudrait-elle si les mille variantes étaient toutes indépendantes~?
\item Quel est le ratio de Sharpe espéré de la meilleure de mille variantes indépendantes sans valeur~?
\end{enumerate}

\medskip\noindent\textbf{Partie III --- Ce que change le journal.}
\begin{enumerate}[resume]
\item Quel est le ratio de Sharpe probabiliste d'un backtest de deux ans à 2.0, lu comme un test unique~?
\item Quel est son ratio de Sharpe déflaté si le journal montre les vingt essais de la semaine~?
\item Et s'il montre les mille de l'année~?
\item Pourquoi mille est-il le bon compte, bien que les variantes soient corrélées~?
\item Combien d'années de backtest faudrait-il pour que la probabilité qu'une variante seule impressionne passe sous 0.001~\%~?
\end{enumerate}

\medskip\noindent\textbf{Partie IV --- Le jalon.}
\begin{enumerate}[resume]
\item La variante impressionnante est envoyée à un \omterm{def:rs:the-research-process:prereg}{échantillon réservé} d'un an, qu'elle réussit si son ratio de Sharpe dépasse 1. Quelle est la
probabilité qu'une variante sans valeur réussisse~?
\item Avec une probabilité a priori de 5~\% et la puissance de l'\omterm{def:rs:the-research-process:prereg}{échantillon réservé} contre un ratio de Sharpe réel de 1.5, quelle est la VPP après l'\omterm{def:rs:the-research-process:prereg}{échantillon réservé}~?
\item Le chercheur a consulté l'\omterm{def:rs:the-research-process:prereg}{échantillon réservé} quatre fois pendant l'ajustement. Quelle est la taille effective de l'\omterm{def:rs:the-research-process:prereg}{échantillon réservé}~?
\item Énoncer le \emph{résultat nommé}~: la probabilité d'au moins un backtest impressionnant en un an de cette recherche, et le
ratio de Sharpe déflaté de ce backtest avec le nombre consigné.
\item Quelles trois choses la présentation du vendredi aurait-elle dû montrer~?
\item En une phrase~: pourquoi le \omterm{def:rs:the-research-process:log}{journal de recherche} est-il un outil de gestion des risques~?
\end{enumerate}
\end{problem}

\section{Questions d'entretien}

\begin{interviewq}[$\star$ \iqroles{chercheur, trader}]\label{iq:rs:the-research-process:1}
Une stratégie candidate affiche un ratio de Sharpe annuel de 2.5 dans son backtest. Quelles sont les trois premières questions que vous posez~?
\end{interviewq}

\begin{interviewq}[$\star\star$ \iqroles{chercheur}]\label{iq:rs:the-research-process:2}
La plupart des idées qu'une équipe teste sont fausses. Pourquoi cela compte-t-il davantage que le niveau de signification de ses tests~?
\end{interviewq}

\begin{interviewq}[$\star\star$ \iqroles{chercheur, ingénieur ML}]\label{iq:rs:the-research-process:3}
Qu'est-ce qui fait qu'un \omterm{def:rs:the-research-process:prereg}{échantillon réservé} cesse d'être un \omterm{def:rs:the-research-process:prereg}{échantillon réservé}~?
\end{interviewq}

\begin{interviewq}[$\star\star$ \iqroles{trader, chercheur, risque}]\label{iq:rs:the-research-process:4}
Une stratégie dont le ratio de Sharpe backtesté est de 1.5 est plate après une année en direct. La arrêtez-vous~?
\end{interviewq}

\begin{interviewq}[$\star\star$ \iqroles{développeur, chercheur}]\label{iq:rs:the-research-process:5}
Concevez un \omterm{def:rs:the-research-process:log}{journal de recherche}. Que contient chaque entrée, et comment le rendez-vous infalsifiable~?
\end{interviewq}

\begin{interviewq}[$\star\star\star$ \iqroles{chercheur}]\label{iq:rs:the-research-process:6}
Donnez une \omterm{def:rs:the-research-process:hypothesis}{justification économique} du renversement de court terme sur les actions, et trois prédictions qu'elle fait en plus du renversement lui-même.
\end{interviewq}
